\section{Sirt taranglik koefitsienti}


\begin{enumerate}

\item Sirt taranglik koeffitsiyentining birligini ko‘rsating.
\begin{enumerate}[label=\Alph*)]
\item $N$
\item $J$
\item $J/m$
\item $J/m^2$
\end{enumerate}

\item Sirt taranglik koeffitsiyenti qanday birlikda ifodalanadi?
\begin{enumerate}[label=\Alph*)]
\item $J \cdot s$
\item $J/m^3$
\item $J/m^2$
\item $N/m$
\end{enumerate}

\item Qo‘zg‘aluvchan simning uzunligi 5 sm bo‘lgan simli ramkadagi sovun pardasi 10 sm ga cho‘zilgan bo‘lsa, cho‘zuvchi kuchning modulini toping (mN). Sovun eritmasining sirt taranglik koeffitsiyenti 40 mN/m.
\begin{enumerate}[label=\Alph*)]
\item 20
\item 2
\item 4
\item 8
\end{enumerate}

\item Diametri 4 mm bo‘lgan kapillyar naychadan tomayotgan suv tomchisining massasini toping (mg). ($\sigma = 73 \, mN/m$).
\begin{enumerate}[label=\Alph*)]
\item 94
\item 187
\item 24
\item 31,4
\end{enumerate}

\item Diametri 1 mm bo‘lgan vertikal shisha naychadan suv tomchisi uzildi. Tomchining og‘irligi qanday (mN)? Suvning sirt taranglik koeffitsiyenti $\sigma = 73 \, mN/m$.
\begin{enumerate}[label=\Alph*)]
\item 0,11
\item 0,23
\item 0,32
\item 0,41
\end{enumerate}

\item Suv diametri 1,8 mm bo‘lgan naychadan tomchilayotgan bo‘lsa, $1 \, cm^3$ suvda nechta tomchi bo‘ladi? Suvning sirt taranglik koeffitsiyenti $\sigma = 72 \, mN/m$.
\begin{enumerate}[label=\Alph*)]
\item 30
\item 25
\item 12
\item 49
\end{enumerate}

\item Teshigining diametri 2,0 mm bo‘lgan tomizgichdan har bir sekundda bittadan kerosin tomchilari tushmoqda. $25 \, cm^3$ kerosin qancha vaqtda (min) oqib tushadi? $\rho = 800 \, kg/m^3$, $\sigma = 24 \, mN/m$.
\begin{enumerate}[label=\Alph*)]
\item 28
\item 16
\item 22
\item 48
\end{enumerate}

\item Radiusi yarim millimetr bo‘lgan tomizgichdan sirt taranglik koeffitsiyenti 50 mN/m bo‘lgan 100 g massali suyuqlik nechta tomchi bo‘lib tomib ketadi?
\begin{enumerate}[label=\Alph*)]
\item 6248
\item 7459
\item 3456
\item 4560
\end{enumerate}

\item Pastga $4,5 \, m/s^2$ tezlanish bilan harakatlanayotgan diametri 1 mm bo‘lgan tomizgichdan uzilgan suv tomchisining massasini toping (mg). ($\sigma = 73 \, mN/m$).
\begin{enumerate}[label=\Alph*)]
\item 18
\item 43,24
\item 24
\item 22,94
\end{enumerate}

\item Hajmi $6 \, cm^3$ bo‘lgan vannachaga naychadan suv tomchilab tushmoqda. Naychaning ichki diametri 1,0 mm ga teng. Vannachani $20^\circ C$ temperaturali suv bilan to‘ldirish uchun naychadan necha tomchi suv tushishi kerak? $\rho = 1000 \, kg/m^3$, $\sigma = 72 \, mN/m$.
\begin{enumerate}[label=\Alph*)]
\item 242
\item 196
\item 225
\item 265
\end{enumerate}

\item Teshigining diametri 2 mm bo‘lgan tomizgichda $219 \, cm^3$ suv bor. Tomizgichdan hammasi bo‘lib nechta tomchi tomadi? $\pi = 3$, $g = 10 \, m/s^2$ deb hisoblang.
\begin{enumerate}[label=\Alph*)]
\item 2180
\item 4000
\item 5000
\item 6000
\end{enumerate}

\item Tomizgich teshigining diametri 1,2 mm. Tomizgichdan oqib chiqayotgan suv tomchisining uzilish momentidagi massasi qancha (mg)?
\begin{enumerate}[label=\Alph*)]
\item 56
\item 28
\item 2,8
\item 5,6
\end{enumerate}

\item Suv ostidagi 18 mm diametrli havo pufagi sirti hosil qiladigan qo‘shimcha bosim qiymatini (Pa) aniqlang. $\sigma = 72 \, mN/m$.
\begin{enumerate}[label=\Alph*)]
\item 32
\item 16
\item 48
\item 34
\end{enumerate}

\item 20 mm diametrli sovun pufagi sirti hosil qiladigan qo‘shimcha bosim qiymatini aniqlang (Pa). $\sigma = 0,04 \, N/m$.
\begin{enumerate}[label=\Alph*)]
\item 12
\item 16
\item 28
\item 40
\end{enumerate}

\item 50 mm diametrli sovun pufagi sirtining ortiqcha potensial energiyasini hisoblang (J). $\sigma = 0,04 \, N/m$.
\begin{enumerate}[label=\Alph*)]
\item $2,4 \cdot 10^{-2}$
\item $24 \cdot 10^{-5}$
\item $8,3 \cdot 10^{-4}$
\item $6,3 \cdot 10^{-4}$
\end{enumerate}

\item Radiusi 4 sm bo‘lgan sovunli suv pufagining sirt energiyasi ($mJ$) nimaga teng? $\sigma = 40 \, mN/m, \pi = 3$.
\begin{enumerate}[label=\Alph*)]
\item 0,77
\item 2,24
\item 1,54
\item 2,3
\end{enumerate}

\item Sovun pufakchasining sirt energiyasi 64 marta ortishi uchun uning radiusi necha marta orttirilishi kerak?
\begin{enumerate}[label=\Alph*)]
\item 16 marta
\item 64 marta
\item 8 marta
\item 4 marta
\end{enumerate}

\item Sovun pufagining radiusini 1 sm dan 6 sm gacha puflab shishirish uchun qanday ish bajarish kerak (mJ)? $\sigma = 45 \, mN/m$.
\begin{enumerate}[label=\Alph*)]
\item 4,5
\item 4
\item 3,6
\item 2
\end{enumerate}

\item Sovun pufagining yuzasini $1 \, cm^2$ dan $3 \, cm^2$ gacha oshirish uchun necha mikrojoul ish bajarish kerak? $\sigma = 5 \cdot 10^{-2} \, N/m$.
\begin{enumerate}[label=\Alph*)]
\item 1
\item 5
\item 20
\item 50
\end{enumerate}

\item Radiusi 0,6 mm bo‘lgan suv tomchisi silliq sirtga urilib 8 ta bir xil kichik tomchiga ajralib ketdi. Bunda tomchi tushgan balandlikning eng kichik qiymati necha sm bo‘ladi? $\sigma = 72 \, mN/m$.
\begin{enumerate}[label=\Alph*)]
\item 2,8
\item 2,6
\item 3,6
\item 4,3
\end{enumerate}

\item Radiusi 0,3 mm bo‘lgan suv tomchisi silliq sirtga urilib 8 ta bir xil kichik tomchiga ajralib ketdi. Bunda tomchi tushgan balandlikning eng kichik qiymati necha sm bo‘ladi? $\sigma = 72 \, mN/m$.
\begin{enumerate}[label=\Alph*)]
\item 4,3
\item 5,2
\item 7,2
\item 8,6
\end{enumerate}

\item Shar shaklidagi sakkizta bir xil haroratli simob tomchilari birlashib yagona tomchi hosil qildi. Bunda harorat qanday va nima sababdan o‘zgaradi?
\begin{enumerate}[label=\Alph*)]
\item Ortadi, chunki sirt energiyasi ortadi.
\item Pasayadi, chunki sirt energiyasi kamayadi.
\item Ko‘tariladi, chunki sirt energiyasi kamayadi.
\item O‘zgarmaydi, chunki sirt energiyasi o'zgarmaydi.
\end{enumerate}

\end{enumerate}



\begin{table}[h!]
\centering
\begin{tabular}{|c|c|c|c|c|c|c|c|c|c|}
\hline
№ & J. & № & J. & № & J. & № & J. & № & J. \\ \hline
1 & D & 6 & B & 11 & C & 16 & C & 21 & C \\ \hline
2 & D & 7 & C & 12 & B & 17 & C & 22 & C \\ \hline
3 & C & 8 & A & 13 & B & 18 & B & & \\ \hline
4 & A & 9 & B & 14 & B & 19 & C & & \\ \hline
5 & B & 10 & D & 15 & B & 20 & C & & \\ \hline
\end{tabular}
\end{table}